% Источники: exam/materials/преподаватель/2020/23-03-2020-MODEL L5(1).doc; exam/materials/моделирование.pdf, раздел 33.
\section{Подходы к оцениванию параметров зависимости между ненаблюдаемыми одновременно характеристиками объекта. Метод максимального правдоподобия.}

\subsection*{Метод максимального правдоподобия (ММП)}

\textbf{Идея:} выбрать значение параметра $\boldsymbol{\theta}$, при котором наблюдаемая выборка $x_1, \ldots, x_n$ наиболее вероятна.

\textbf{Функция правдоподобия} (для непрерывного распределения):
\[
  L(x_1,\ldots,x_n;\,\boldsymbol{\theta}) = \prod_{i=1}^{n} f(x_i;\,\boldsymbol{\theta}).
\]

\textbf{Оценка ММП:}
\[
  \hat{\boldsymbol{\theta}}_{\text{ММП}} = \arg\max_{\boldsymbol{\theta}} L(x_1,\ldots,x_n;\,\boldsymbol{\theta}).
\]

На практике максимизируют \textbf{логарифмическую функцию правдоподобия}:
\[
  \ell(\boldsymbol{\theta}) = \ln L = \sum_{i=1}^n \ln f(x_i;\,\boldsymbol{\theta}),
\]
решая \textbf{уравнение правдоподобия}:
\[
  \frac{\partial \ell}{\partial \boldsymbol{\theta}} = \mathbf{0}.
\]

\subsection*{Свойства ММП-оценки}

\begin{itemize}
  \item Когда существует эффективная оценка, метод максимального правдоподобия даёт эту оценку.
  \item В общем случае при достаточно общих условиях метод максимального правдоподобия даёт асимптотически несмещённые и асимптотически эффективные оценки при $n\to\infty$.
\end{itemize}

\subsection*{Пример: логистическая регрессия}

Логистическая регрессия — это статистическая модель, используемая для определения вероятности принадлежности объекта $X_i$ к одному из двух классов. Параметры такой модели выбираются по тому же принципу: берутся значения, при которых функция правдоподобия наблюдаемой выборки обращается в максимум.

\clearpage
